% Created 2026-03-16 Mon 17:28
% Intended LaTeX compiler: lualatex
\documentclass[11pt]{article}
\usepackage{amsmath}
\usepackage{fontspec}
\usepackage{graphicx}
\usepackage{longtable}
\usepackage{wrapfig}
\usepackage{rotating}
\usepackage[normalem]{ulem}
\usepackage{capt-of}
\usepackage{hyperref}
\author{Prologos Project}
\date{2026-03-16}
\title{Quantum-Like Models for Social, Economic, and Cognitive Systems\\\medskip
\large A Research Background for Tensor-Aware Propagator Networks in Prologos}
\hypersetup{
 pdfauthor={Prologos Project},
 pdftitle={Quantum-Like Models for Social, Economic, and Cognitive Systems},
 pdfkeywords={},
 pdfsubject={},
 pdfcreator={Emacs 30.2 (Org mode 9.7.39)}, 
 pdflang={English}}
\begin{document}

\maketitle
\tableofcontents

\section{1. Introduction}
\label{sec:org1f81317}

\subsection{1.1 The Quantum Probability Revolution in Social Science}
\label{sec:org676a9dc}

Over the past three decades, a growing community of researchers has discovered
that the mathematical formalism of quantum mechanics --- Hilbert spaces,
superposition, interference, entanglement, and projective measurement --- provides
a more natural framework than classical (Kolmogorovian) probability for modeling
human decision-making, cognition, social dynamics, and economic behavior. This is
not a claim about quantum physics in the brain. It is a claim about mathematics:
quantum probability theory is the \emph{next-simplest} probability theory after
classical probability, and it naturally accommodates phenomena --- contextuality,
order effects, interference --- that classical probability cannot represent.

The field has reached institutional maturity. A dedicated journal, \emph{Quantum
Economics and Finance} (SAGE, in collaboration with the Center for Quantum Social
and Cognitive Science), publishes peer-reviewed research. Special issues in
\emph{Quantum Reports} (MDPI), \emph{Frontiers in Physics}, and \emph{Topics in Cognitive
Science} have collected major contributions. Comprehensive references include
Busemeyer \& Bruza's \emph{Quantum Models of Cognition and Decision} (2012), Haven \&
Khrennikov's \emph{Quantum Social Science} (2013), and the \emph{Palgrave Handbook of
Quantum Models in Social Science}. A 2026 systematic review (Jammal et al.)
documents steady publication growth, with over 60 articles in 2024 alone.
\subsection{1.2 Why Hilbert Space?}
\label{sec:org1768d98}

Classical probability (Kolmogorov, 1933) assigns probabilities to events as
elements of a Boolean algebra. The key properties: probabilities are real and
non-negative; they add for disjoint events; the total is 1. This framework
assumes:
\begin{itemize}
\item \textbf{Commutativity:} P(A and B) = P(B and A).
\item \textbf{The law of total probability:} P(A) = P(A|B)P(B) + P(A|¬B)P(¬B).
\item \textbf{Single sample space:} All events can be jointly represented.
\end{itemize}

Human cognition and social behavior routinely violate all three. The order in
which questions are asked changes the answers (non-commutativity). The disjunction
effect violates the law of total probability. Preferences cannot always be
represented in a single consistent framework (contextuality).

Quantum probability (von Neumann, 1932) replaces the Boolean algebra with a
Hilbert space. Events are subspaces; probabilities are squared magnitudes of
projections. This framework accommodates:
\begin{itemize}
\item \textbf{Non-commutativity:} Projecting onto subspace A then B differs from B then A.
\item \textbf{Interference:} |a+b|² ≠ |a|² + |b|² when cross-terms are nonzero.
\item \textbf{Contextuality:} Different measurement contexts yield incompatible probability
spaces.
\end{itemize}

The crucial point: these are \emph{mathematical} properties of Hilbert spaces, not
\emph{physical} properties of quantum mechanics. Any system exhibiting contextuality,
order effects, or interference --- whether it is a photon, a human decision-maker,
or a financial market --- is more naturally described by quantum probability than
by classical probability.
\subsection{1.3 Historical Overview}
\label{sec:org9753fb5}

The idea of applying quantum formalism to cognition emerged independently from
several sources in the 1990s:

\begin{itemize}
\item \textbf{Aerts and collaborators} (Brussels) developed a quantum-like model of concept
combination, showing that the "pet-fish" problem (a guppy is a good example of
"pet fish" but a poor example of either "pet" or "fish") requires a state space
richer than classical probability.
\item \textbf{Khrennikov} (Linnaeus University) developed a contextual probability framework,
pointing out that violations of the formula of total probability can serve as
a quantitative measure of non-classicality.
\item \textbf{Bordley} applied quantum probability to decision-making, showing that
interference effects explain preference reversals.
\end{itemize}

The field accelerated in the 2000s with:
\begin{itemize}
\item \textbf{Busemeyer and collaborators} (Indiana University) developing formal quantum
models of specific cognitive phenomena: the conjunction fallacy, the disjunction
effect, order effects in judgments.
\item \textbf{Haven and Khrennikov} connecting quantum formalism to financial markets and
economics.
\item \textbf{Yukalov and Sornette} (ETH Zurich) developing Quantum Decision Theory as a
comprehensive mathematical framework.
\end{itemize}

The 2010s saw consolidation into textbooks and institutional infrastructure.
The 2020s have brought empirical validation at scale (the QQ equality tested on
70+ national surveys), practical applications (quantum computing for portfolio
optimization), and new theoretical developments (quantum-like Bayesian networks,
open quantum systems for cognitive dynamics).
\subsection{1.4 Connection to Prologos}
\label{sec:orgd574b2f}

A companion document ("Tensor-Aware Propagator Networks with Unitary Evolution")
describes how Prologos's propagator network can be extended with TensorCells
(joint amplitude distributions over multiple variables) and unitary evolution
(reversible transformations of quantum-like states). This background document
provides the domain knowledge for what such a system could model.

The connection is direct: every quantum-like model in the literature requires
(1) a state space (Hilbert space), (2) evolution operators (unitaries), and
(3) measurement (projection + Born rule). In Prologos, these map to
(1) TensorCells, (2) unitary propagators, and (3) the measurement stratum with
ATMS integration. The dependent type system provides compile-time dimensional
correctness; the propagator network provides the computational engine; the
Layered Recovery Principle ensures that unitary evolution and measurement are
correctly stratified.
\section{2. Foundations: Quantum Probability vs. Classical Probability}
\label{sec:org49de79e}

\subsection{2.1 Kolmogorov Probability and Its Limitations}
\label{sec:org6f10e76}

Classical probability theory (Kolmogorov, 1933) rests on a probability space
(Ω, F, P) where Ω is a sample space, F is a σ-algebra of events, and P is a
probability measure. For finite spaces, this reduces to:
\begin{itemize}
\item A set of outcomes \{ω₁, \ldots{}, ωₙ\}.
\item Probabilities p(ωᵢ) ≥ 0 with Σᵢ p(ωᵢ) = 1.
\item P(A) = Σ\textsubscript{ωᵢ∈A} p(ωᵢ) for any event A ⊆ Ω.
\end{itemize}

The law of total probability follows: P(A) = P(A|B)P(B) + P(A|B\textsuperscript{c})P(B\textsuperscript{c}),
expressing P(A) as a weighted sum over a partition. This is the foundation of
Bayesian reasoning.

\textbf{Limitations in cognitive science:}

The law of total probability is empirically violated in human judgment.
Tversky and Shafir (1992) demonstrated the \textbf{disjunction effect}: in a
two-stage gamble, participants who knew they won the first stage chose to play
again (67\%); participants who knew they lost also chose to play again (59\%);
but participants who did \emph{not} know the outcome of the first stage chose to play
again at a lower rate (36\%). Under classical probability, not knowing the
outcome should produce a probability between 59\% and 67\%, not below both.
This is a direct violation of the law of total probability.
\subsection{2.2 Hilbert Space Probability}
\label{sec:org77c741b}

Quantum probability replaces the classical framework with:
\begin{itemize}
\item A Hilbert space H (a complex vector space with inner product).
\item A state |ψ⟩ ∈ H with ||ψ|| = 1 (a unit vector).
\item Observables as projection operators P\textsubscript{A} (self-adjoint, P² = P).
\item Probability of event A: p(A) = ||P\textsubscript{A}|ψ⟩||² = ⟨ψ|P\textsubscript{A}|ψ⟩.
\item State update after observing A: |ψ'⟩ = P\textsubscript{A}|ψ⟩ / ||P\textsubscript{A}|ψ⟩|| (Lüders' rule).
\end{itemize}

\textbf{Key differences from classical probability:}

\begin{enumerate}
\item \textbf{Non-commutativity.} If P\textsubscript{A} P\textsubscript{B} ≠ P\textsubscript{B} P\textsubscript{A}, then the probability of "A then B"
differs from "B then A":
p(A then B) = ||P\textsubscript{B} P\textsubscript{A}|ψ⟩||² ≠ ||P\textsubscript{A} P\textsubscript{B}|ψ⟩||² = p(B then A).
This models order effects in survey questions.

\item \textbf{Interference.} If |ψ⟩ = α|a⟩ + β|b⟩, then:
p(c) = |⟨c|ψ⟩|² = |α⟨c|a⟩ + β⟨c|b⟩|²
     = |α|²|⟨c|a⟩|² + |β|²|⟨c|b⟩|² + 2Re(α*β⟨c|a⟩*⟨c|b⟩)
The last term is the \emph{interference term}, absent in classical probability.

\item \textbf{Incompatibility.} Two observables A, B are incompatible if their projection
operators do not commute. Measuring A disturbs the state with respect to B.
There is no joint probability distribution for incompatible observables.
\end{enumerate}
\subsection{2.3 Interference and Superposition}
\label{sec:org8ebe4f4}

The interference term is the key mechanism by which quantum-like models produce
non-classical predictions. Consider a decision between options C (cooperate) and
D (defect):

\begin{itemize}
\item Classical: p(C) = p(C|know\textsubscript{win})p(know\textsubscript{win}) + p(C|know\textsubscript{lose})p(know\textsubscript{lose})
= 0.67 × 0.5 + 0.59 × 0.5 = 0.63.
\item Quantum: |ψ⟩ = α|win⟩ + β|lose⟩. After interference:
p(C) = |α|²p(C|win) + |β|²p(C|lose) + 2Re(α*β × interference)
The interference term can be \emph{negative}, reducing p(C) below both conditional
probabilities. This explains the disjunction effect.
\end{itemize}

The double-slit analogy is useful: when a particle passes through two slits
simultaneously (superposition), the probability pattern at the detector shows
interference fringes. When a decision-maker considers two possible states of
the world simultaneously (uncertainty), the choice probability shows
"interference fringes" --- deviations from the weighted average of the
conditional probabilities.
\subsection{2.4 Entanglement and Non-Separability}
\label{sec:org35b767f}

Two systems A and B are \textbf{entangled} when their joint state cannot be factored:
\begin{center}
\begin{tabular}{lllll}
ψ\textsubscript{AB}⟩ ≠ & ψ\textsubscript{A}⟩ ⊗ & ψ\textsubscript{B}⟩ for any states & ψ\textsubscript{A}⟩, & ψ\textsubscript{B}⟩.\\
\end{tabular}
\end{center}

The Bell state (|00⟩ + |11⟩)/√2 is the canonical example: if A is measured and
found to be 0, then B is guaranteed to be 0 as well, even though neither had a
definite value before measurement. The correlations are stronger than any
classical probabilistic model can produce (Bell's theorem).

In social science, entanglement models non-separable preferences and correlated
beliefs:
\begin{itemize}
\item \textbf{Voter preferences:} A voter's preference for president and senator cannot always
be factored into independent presidential and senatorial preferences. "Ticket
splitting" (different parties for different races) may violate classical
separability.
\item \textbf{Agent interactions:} Two traders whose strategies are correlated through shared
information or social influence may have entangled strategy states.
\item \textbf{Concept combination:} The concept "pet fish" (well-modeled by a guppy) cannot
be derived from independent concepts "pet" (well-modeled by a dog) and "fish"
(well-modeled by a trout). Aerts showed this requires entangled concept states.
\end{itemize}
\subsection{2.5 Open Quantum Systems and Decoherence}
\label{sec:orgf5dba0c}

Real quantum systems interact with their environments, causing decoherence --- the
loss of quantum coherence (interference capability). The GKSL (Gorini-Kossakowski-
Sudarshan-Lindblad) master equation governs this evolution:

\begin{verbatim}
dρ/dt = -i[H, ρ] + Σ_k (L_k ρ L_k† - ½{L_k†L_k, ρ})
\end{verbatim}

where ρ is the density matrix (mixed state), H is the Hamiltonian (unitary
dynamics), and L\textsubscript{k} are Lindblad operators (environmental interaction). The first
term drives unitary evolution (rotation); the second term drives decoherence
(decay of off-diagonal elements).

In cognitive modeling (Asano, Khrennikov et al.), this models decision-making:
\begin{itemize}
\item The initial state ρ₀ is a superposition of choices (indecision).
\item The Hamiltonian H encodes the decision-maker's preferences and information.
\item The Lindblad operators L\textsubscript{k} model interaction with the "mental environment"
(memory, emotions, social pressure).
\item The stationary state ρ\_∞ is the final decision --- a classical mixture
(diagonal density matrix) with no remaining quantum coherence.
\end{itemize}

The \textbf{diffidence} (classical entropy minus quantum entropy) tracks the decision
process: it starts high (quantum uncertainty exceeds classical) and decays to zero
as decoherence eliminates quantum effects, leaving a classical choice.
\subsection{2.6 Contextuality}
\label{sec:orgf099a5b}

\textbf{Contextuality} is the property that the outcome of a measurement depends on what
other measurements are performed alongside it. In classical probability, outcomes
are context-independent: the probability of "heads" on a coin does not depend on
whether a die is also rolled. In quantum mechanics, the probability of a spin
measurement outcome depends on which axis is measured simultaneously.

Dzhafarov and Kujala developed the \textbf{Contextuality-by-Default (CbD)} framework,
which extends contextuality analysis to "inconsistently connected" systems --- those
where marginal distributions change across contexts (a common feature of behavioral
data). The CbD framework distinguishes between:
\begin{itemize}
\item \textbf{Direct influences:} Changes in marginal distributions due to context changes.
\item \textbf{True contextuality:} Correlations that cannot be explained by any classical
coupling, even after accounting for direct influences.
\end{itemize}

The "Snow Queen" experiment (Cervantes \& Dzhafarov) provided the first
unequivocal empirical evidence of true contextuality in human behavior. A 2025
study found true contextuality in global investment decisions, demonstrating that
the phenomenon extends beyond laboratory settings to real-world economic behavior.
\section{3. Quantum Decision Theory}
\label{sec:org7b8b33c}

\subsection{3.1 The Yukalov-Sornette Framework}
\label{sec:org7de854a}

Quantum Decision Theory (QDT), developed by Vyacheslav Yukalov and Didier Sornette
(ETH Zurich), is a comprehensive mathematical framework that decomposes behavioral
probability into classical and quantum components.

\textbf{Mathematical formulation.} The probability of choosing a prospect πⱼ is:

\begin{verbatim}
p(πⱼ) = f(πⱼ) + q(πⱼ)
\end{verbatim}

where:
\begin{itemize}
\item f(πⱼ) is the \textbf{utility factor} --- the classical probability derived from
expected utility or other rational choice models. It satisfies all classical
probability axioms: f(πⱼ) ≥ 0, Σⱼ f(πⱼ) = 1.
\item q(πⱼ) is the \textbf{attraction factor} --- the quantum interference term capturing
emotional, irrational, and bias-driven components. It satisfies Σⱼ q(πⱼ) = 0
(the interference terms sum to zero, preserving total probability).
\end{itemize}

The attraction factor arises from the entanglement between the decision-maker's
strategic state and the prospects being evaluated. When prospects are composite
(involving multiple attributes or stages), the probability operator decomposes
into products of sub-operators, and cross-terms appear due to entanglement.

\textbf{The Quarter Law.} QDT predicts that the average magnitude of the attraction factor
is:

\begin{verbatim}
⟨|q|⟩ = 1/4
\end{verbatim}

This is not merely an empirical finding but is derivable as the non-informative
prior over several families of distributions (uniform, beta, Gaussian). The
Quarter Law provides a parameter-free quantitative prediction: the quantum
deviation from classical expected utility averages 25\% of the choice probability.

\textbf{Empirical validation.} The Quarter Law has been tested on simple risky choices
(lottery pairs) and group-level data strongly support it (Yukalov \& Sornette, 2016).
\subsection{3.2 Classical Decision Paradoxes Resolved}
\label{sec:org0278a82}

QDT resolves the major anomalies of classical decision theory as natural
consequences of quantum interference:

\textbf{Allais paradox.} People prefer a certain \$1M to an 89\% chance of \$1M + 10\%
chance of \$5M + 1\% chance of \$0, yet simultaneously prefer a 10\% chance of \$5M
to an 11\% chance of \$1M. This violates expected utility's independence axiom.
In QDT, the interference term is positive for the certain option (reinforcing
certainty preference) and negative for the mixed gamble (suppressing risk
tolerance), producing the observed pattern.

\textbf{Ellsberg paradox.} People prefer betting on a known 50-50 urn to an urn with
unknown composition, even when the expected values are identical. Classical
probability requires ambiguity-neutral behavior. In QDT, ambiguity generates
a negative attraction factor (ambiguity aversion), reducing the probability
of choosing the ambiguous option below its utility factor.

\textbf{Disjunction effect.} As described in §2.1: knowing the outcome of a first stage
leads to playing again regardless, but not knowing leads to refusing. The
interference between the "won" and "lost" branches is destructive,
reducing the play-again probability below both conditional probabilities.

\textbf{Conjunction fallacy.} People judge P(feminist ∧ bank teller) > P(bank teller),
violating the conjunction rule. In QDT, the conjunction is not a classical
intersection but a sequential projection. The "feminist" projection rotates
the state into a subspace where "bank teller" has higher probability than
it does from the initial state.
\subsection{3.3 Quantum Prospect Theory}
\label{sec:org6e88037}

Quantum Prospect Theory augments Kahneman and Tversky's Cumulative Prospect
Theory (CPT) with quantum interference effects. The key formula adds an
interference term to the CPT value function:

\begin{verbatim}
V_quantum(prospect) = V_CPT(prospect) + interference(prospect)
\end{verbatim}

The interference term captures conflicting interests, ambiguity, and emotions
arising during deliberation. Two estimation methods exist:
\begin{itemize}
\item The \textbf{interference alternation method}: estimates the sign and magnitude of
interference from the structure of the prospect.
\item The \textbf{interference quarter law}: uses the QDT quarter law as a default
magnitude.
\end{itemize}

Empirical comparisons show that at the aggregate level, CPT-based QDT outperforms
pure CPT. At the individual level, RDU-based (Rank-Dependent Utility) QDT best
captures heterogeneity across subjects.
\subsection{3.4 Order Effects in Judgments}
\label{sec:org3ebc165}

Wang and Busemeyer (2013) derived a parameter-free prediction from quantum
probability: the \textbf{QQ equality}.

For two binary yes/no questions A and B:

\begin{verbatim}
p(Ay,By) + p(An,Bn) - p(By,Ay) - p(Bn,An) = 0
\end{verbatim}

where p(Ay,By) is the probability of answering "yes" to A and then "yes" to B.

\textbf{Derivation.} The QQ equality follows from the law of reciprocity for Hilbert
space inner products: |⟨ψ\textsubscript{B}|ψ\textsubscript{A}⟩|² = |⟨ψ\textsubscript{A}|ψ\textsubscript{B}⟩|². The real part of the inner
product is independent of projection order, yielding the equality.

\textbf{Empirical validation.} Wang et al. (2014) tested the QQ equality on \textbf{70 national
representative surveys} (most with >1000 participants each, conducted by Gallup
and Pew Research) plus two laboratory experiments. The QQ equality was strongly
upheld --- a remarkable result given that it is a parameter-free, a priori
prediction derived purely from the mathematical structure of Hilbert space.

Neither Bayesian models nor Markov models generally satisfy the QQ equality.
Its empirical success is one of the strongest pieces of evidence for the quantum
probability framework in cognitive science.
\section{4. Quantum Cognition}
\label{sec:org7e47156}

\subsection{4.1 Similarity Judgments}
\label{sec:orgdca1fda}

The \textbf{Quantum Similarity Model} (QSM), developed by Pothos and Busemeyer, models
concepts as subspaces of a Hilbert space (not as points, as in classical geometric
models). The key innovation: similarity between two concepts is computed as the
probability of a sequential projection.

\textbf{Formal structure:}
\begin{itemize}
\item A concept C is associated with a subspace H\textsubscript{C} of the Hilbert space.
\item The mental state is a vector |ψ⟩.
\item Similarity of A to B is: sim(A,B) = ||P\textsubscript{B} P\textsubscript{A}|ψ⟩||² (project onto A, then B).
\end{itemize}

\textbf{Key advantages:}
\begin{itemize}
\item \textbf{Asymmetry.} sim(A,B) ≠ sim(B,A) because P\textsubscript{B} P\textsubscript{A} ≠ P\textsubscript{A} P\textsubscript{B} in general. This
captures Tversky's (1977) finding that similarity judgments are asymmetric:
North Korea is more similar to China than China is to North Korea.
\item \textbf{Diagnosticity.} Context (prior projections) changes similarity judgments. After
seeing several fruits, an apple is similar to an orange; after seeing several
round objects, an apple is similar to a ball.
\item \textbf{Unification.} The QSM uses the same mathematical framework as the conjunction
fallacy model, the order effects model, and the disjunction effect model. All
are instances of projection in Hilbert space.
\end{itemize}
\subsection{4.2 Conceptual Combination}
\label{sec:org4871c8f}

The conjunction fallacy (§3.2) is a special case of conceptual combination. More
broadly, quantum models address the general problem of how concepts combine:

\textbf{The guppy problem.} A guppy is a good example of "pet fish" but a poor example
of either "pet" or "fish." Classical fuzzy set theory (Osherson \& Smith, 1981)
predicts that membership in A∧B should be at most min(membership in A, membership
in B). The guppy violates this: membership in "pet fish" exceeds membership in
both components.

Aerts (2009) showed that this requires an \emph{entangled} state in the tensor product
of the "pet" and "fish" Hilbert spaces. The concept "pet fish" is not the
intersection of "pet" and "fish" but a projection onto a subspace of the composite
space that has no product-state representation.

\textbf{Overextension.} When combining concepts, people sometimes extend the category
beyond what either component warrants. "Stone lion" is judged as a better
predator than "stone" warrants and a better decorative object than "lion" warrants.
This "emergence of new meaning" in conceptual conjunction maps to the quantum
phenomenon of measurement outcomes that differ from what either subsystem alone
would produce.
\subsection{4.3 Memory and Perception}
\label{sec:org48d7372}

Quantum models have been applied to:
\begin{itemize}
\item \textbf{Memory retrieval:} The order in which memory items are recalled affects
subsequent retrieval probabilities, modeled as sequential projections.
\item \textbf{Perception:} Ambiguous stimuli (Necker cube, duck-rabbit) are modeled as
superpositions that collapse upon "measurement" (conscious perception).
The dynamics of perceptual switching are modeled by unitary evolution.
\item \textbf{Categorization:} Categorization decisions under uncertainty exhibit
interference effects analogous to the disjunction effect.
\end{itemize}
\subsection{4.4 Causal Reasoning}
\label{sec:org4b56273}

Trueblood and Busemeyer (2011, 2012) developed \textbf{quantum-like Bayesian networks}
(QLBNs) by replacing classical probabilities with quantum probability amplitudes.

\textbf{Framework:}
\begin{itemize}
\item Any classical Bayesian network can be generalized to a quantum Bayesian network
by replacing conditional probability tables with conditional amplitude tables.
\item A \textbf{similarity heuristic} automatically fits parameters through vector
similarities, addressing the exponential growth of quantum parameters.
\end{itemize}

\textbf{Key results:}
\begin{itemize}
\item QLBNs account for \textbf{order effects in causal judgments}: judging whether A causes
C and then whether B causes C gives different results than the reverse order.
Classical Bayesian networks, which satisfy the local Markov condition, cannot
produce order effects.
\item \textbf{Anti-discounting:} In classical causal reasoning, learning about an alternative
cause should reduce belief in the original cause (discounting). Humans sometimes
do the opposite (anti-discounting). QLBNs model this via constructive
interference between causal pathways.
\item \textbf{Reciprocity:} Observing that A causes B increases belief that B causes A, even
when there is no causal mechanism. This arises naturally from the symmetry of
Hilbert space inner products.
\end{itemize}
\subsection{4.5 The Open Systems Approach}
\label{sec:orga5b3af6}

Asano, Ohya, Tanaka, Basieva, and Khrennikov developed a comprehensive model
of cognitive dynamics using the GKSL master equation (§2.5).

\textbf{Core model:}
\begin{enumerate}
\item The decision-maker's mental state is a density matrix ρ on a Hilbert space
whose basis elements correspond to choices.
\item The Hamiltonian H encodes preferences and information dynamics.
\item Lindblad operators L\textsubscript{k} model interaction with the mental environment:
memory retrieval, emotional processing, social pressure.
\item The master equation dρ/dt = -i[H,ρ] + Σ\textsubscript{k}(L\textsubscript{k} ρ L\textsubscript{k}† - ½\{L\textsubscript{k}†L\textsubscript{k}, ρ\})
drives the mental state toward a stationary state ρ\_∞.
\item The stationary state is a classical mixture (diagonal) representing the
final decision.
\end{enumerate}

\textbf{Applications:}
\begin{itemize}
\item \textbf{Disjunction effect in Prisoner's Dilemma:} The Shafir-Tversky experiments
are explained by decoherence: knowing the opponent's move collapses the
superposition, while not knowing preserves it, and interference between
the "opponent cooperated" and "opponent defected" branches reduces the
cooperation probability.
\item \textbf{US political decision-making:} Voters' superposition of Republican/Democrat
preferences decoheres through exposure to campaign information.
\item \textbf{Biological applications:} The same formalism models E. coli glucose-lactose
metabolism and epigenetic evolution.
\end{itemize}

\textbf{Non-Markovian extensions} are an active frontier. The GKSL equation assumes
Markovian (memoryless) dynamics, but human cognition has strong memory effects.
Non-Markovian quantum dynamics (Nakajima-Zwanzig equations, time-convolutionless
approaches) are under development.
\section{5. Quantum Game Theory}
\label{sec:org8b8b78e}

\subsection{5.1 The EWL Protocol}
\label{sec:orgca6e979}

Eisert, Wilkens, and Lewenstein (1999) proposed the first formal framework for
quantum game theory, using entangled quantum strategies.

\textbf{Mathematical formulation:}

\begin{enumerate}
\item Classical strategies are mapped to qubit states: Cooperate = |0⟩, Defect = |1⟩.

\item An \textbf{entangling gate} creates initial correlations:
J(γ) = exp(iγ σ\textsubscript{x} ⊗ σ\textsubscript{x} / 2)
where γ ∈ [0, π/2] controls the entanglement degree and σ\textsubscript{x} is the Pauli-X
matrix.

\item The initial state is: |ψ\textsubscript{in}⟩ = J(γ)|00⟩.

\item Each player applies a unitary strategy U ∈ SU(2), parameterized as:
U(θ, φ) = [[e\textsuperscript{iφ}cos(θ/2), sin(θ/2)], [-sin(θ/2), e\textsuperscript{-iφ}cos(θ/2)]]
With I = U(0,0) = Cooperate and σ\textsubscript{x} = U(π,0) = Defect as classical strategies.

\item The final state is: |ψ\textsubscript{out}⟩ = J†(γ) · (U\textsubscript{A} ⊗ U\textsubscript{B}) · J(γ)|00⟩.

\item Payoffs are computed from measurement probabilities:
P\textsubscript{σ,τ} = |⟨σ,τ|ψ\textsubscript{out}⟩|² for σ,τ ∈ \{0,1\}
Expected payoff = Σ\textsubscript{σ,τ} P\textsubscript{σ,τ} × payoff(σ,τ).
\end{enumerate}

\textbf{Key special strategy:} Q = U(0, π/2) --- a quantum strategy with no classical
analogue. It commutes with J(γ) and has the property that Q against any classical
strategy yields at least the cooperative payoff.
\subsection{5.2 Quantum Prisoner's Dilemma}
\label{sec:orga7c6757}

In the classical Prisoner's Dilemma, the unique Nash equilibrium (Defect, Defect)
is Pareto-dominated by (Cooperate, Cooperate). Rational self-interest leads to a
collectively worse outcome.

The quantum version resolves this:

\textbf{Three entanglement regimes:}
\begin{itemize}
\item γ = 0 (no entanglement): Classical game recovered exactly. J(0) = I, so the
entangling gate is trivial. The unique NE is (D, D).
\item 0 < γ < π/2 (partial entanglement): Intermediate regime. The classical NE
may still exist but with modified payoffs.
\item γ = π/2 (maximum entanglement): The strategy pair (Q, Q) is a Nash equilibrium
\emph{and} Pareto-optimal. Neither player can improve their payoff by unilateral
deviation. The dilemma is resolved.
\end{itemize}

\textbf{The mechanism:} At maximum entanglement, the initial state J(π/2)|00⟩ creates
quantum correlations between the players. The quantum strategy Q exploits these
correlations: Q ⊗ Q against the entangled initial state produces the cooperative
outcome with certainty. Any unilateral deviation from Q reduces the deviator's
payoff (Nash equilibrium property) while no outcome Pareto-dominates Q ⊗ Q.

\textbf{Critiques:} Benjamin and Hayden noted that the two-parameter restriction
U(θ,φ) is not the most general SU(2) strategy. With the full three-parameter
SU(2), Nash equilibria may not exist. The robustness of quantum game-theoretic
advantages under noise and decoherence is an active research area.
\subsection{5.3 Quantum Auctions and Mechanism Design}
\label{sec:orgacd0062}

Quantum information theory has implications for economic mechanism design:

\textbf{Circumventing impossibility results.} Classical mechanism design is constrained
by impossibility theorems (Arrow, Gibbard-Satterthwaite). Agents with access to
quantum mechanisms can, under certain conditions, circumvent these constraints:
quantum entanglement provides correlation resources unavailable classically.

\textbf{Quantum economies.} Players producing and consuming "quantum goods" with
entangled strategies yield economic equilibria with distinctive properties.
These models "shed new light on theories of mechanism design, auction and
contract in the quantum era."

\textbf{Practical quantum auctions.} Economic applications relying on quantum
information aspects (not computational speedup) are viable with just a few
qubits. This suggests potential early applications of near-term quantum
technology in auction design.

\textbf{Quantum Condorcet Voting.} Quantum Condorcet Voting has been used to disprove
Arrow's Impossibility Theorem in the quantum setting: for small readout
amplitudes, quantum outcomes track the classical scheme, but beyond a modest
noise level, they converge to non-classical distributions that satisfy the
conditions Arrow showed to be classically impossible.
\subsection{5.4 Evolutionary Quantum Game Theory}
\label{sec:org084932c}

Quantum game theory extends to evolutionary settings where strategies evolve
over time under selection pressure:

\begin{itemize}
\item \textbf{Quantum replicator dynamics:} The classical replicator equation is generalized
to track evolution of quantum strategy parameters. Quantum strategies can
invade classical populations under certain conditions.
\item \textbf{Biological applications:} Quantum game models have been applied to evolutionary
biology, where the "entanglement" is reinterpreted as genetic correlation
between related organisms.
\item A 2025 review in \emph{Quantum Information Processing} maps the intersection of
quantum and evolutionary game theory comprehensively.
\end{itemize}
\section{6. Quantum Finance and Economics}
\label{sec:org1525abd}

\subsection{6.1 The Schrödinger-Black-Scholes Connection}
\label{sec:org164ac0d}

Emmanuel Haven demonstrated that the Black-Scholes-Merton equation for option
pricing is a special case of the Schrödinger equation when markets are perfectly
efficient. The derivation introduces a parameter ℏ (by analogy with Planck's
constant) that measures the degree of market inefficiency --- the amount of
arbitrage arising from non-instantaneous price changes, information asymmetry,
and unequal trader wealth.

When ℏ → 0 (perfect efficiency), the equation reduces to classical Black-Scholes.
When ℏ > 0, quantum-like effects emerge: interference between trading strategies,
non-commutative observables for price and momentum, and wave-like propagation of
market information. The Black-Scholes equation is equivalent to the imaginary-time
Schrödinger equation of a free particle --- a connection that brings the full
apparatus of quantum mechanics into financial modeling.

The "financial Hamiltonian" (analogous to the energy operator in physics) is
determined by the rate of return. Born's rule, interference of probabilities, and
non-commutative operators all emerge from the organization of production and
investment processes.
\subsection{6.2 Open Quantum Systems in Markets}
\label{sec:orge6ea420}

Market dynamics are naturally modeled as \emph{open} quantum systems --- systems
interacting with an environment. The GKSL master equation (§2.5) governs the
evolution of market states:

\begin{itemize}
\item The density matrix ρ represents the distribution over market configurations.
\item The Hamiltonian H encodes fundamental valuations and market microstructure.
\item Lindblad operators L\textsubscript{k} model information arrival, regulatory intervention,
and trader behavioral noise.
\item Decoherence (decay of off-diagonal elements) corresponds to the market
"resolving" uncertainty --- the transition from quantum-like superposition of
prices to a definite traded price.
\end{itemize}

Khrennikov developed a contextual probability framework using agents to model
this: each trader operates in a context (their information set, risk tolerance,
time horizon), and the incompatibility of different agents' contexts produces
non-classical aggregate behavior.
\subsection{6.3 Portfolio Optimization}
\label{sec:orgc743081}

Quantum computing is finding practical application in portfolio optimization:

\textbf{IBM + Vanguard (2024-2025):} Used 109 qubits on the IBM Quantum Heron r1
processor (133 available) with circuits containing up to 4,200 gates. On
simplified ETF portfolio construction, the hybrid quantum-classical workflow
performed on par with and sometimes surpassed state-of-the-art classical solvers
as problem size increased.

\textbf{Citi + Classiq:} Applied the Quantum Approximate Optimization Algorithm (QAOA)
to portfolio optimization problems.

\textbf{Fidelity + IonQ:} Developed quantum models for synthetic financial data
generation.

Classical portfolio optimization becomes intractable when adding realistic
constraints (lot sizes, maximum positions, sector limits). Quantum computers
encode exponentially many portfolio configurations as superpositions, exploring
the solution space more efficiently and potentially avoiding local minima.

McKinsey estimates \$400-600 billion in potential value from quantum computing in
financial services by 2035. The CFA Institute (2025) recommends starting with
Monte Carlo methods for pricing/risk, then moving to discrete optimization as
quantum hardware matures.
\subsection{6.4 Quantum-Like Models of Market Behavior}
\label{sec:org97a8c1b}

Beyond formal quantum computing, quantum-like models explain market phenomena
that classical models struggle with:

\begin{itemize}
\item \textbf{Interference in trading decisions:} A trader considering two information sources
simultaneously may make different decisions than if they processed each source
sequentially. The interference term captures "cognitive momentum" in trading.
\item \textbf{Non-commutative observables:} Asking "what is the fair price?" and "what is the
market sentiment?" in different orders yields different answers --- the act of
assessing price changes one's sentiment assessment, and vice versa.
\item \textbf{Contextuality in markets:} The same asset has different risk profiles depending
on what else is in the portfolio --- a form of quantum contextuality.
\end{itemize}
\section{7. Quantum Social Dynamics}
\label{sec:org2e4f451}

\subsection{7.1 Opinion Dynamics and Polarization}
\label{sec:orgbb460df}

Recent models (2024-2025) apply quantum formalism to opinion dynamics:

\textbf{Beliefs as quantized energy levels.} A 2024 model represents political beliefs
as energy levels in a quantum system. Polarization corresponds to the formation
of energy bands and interband gaps when like-minded individuals interact.
Transitions between polarized groups require significant "energy" (effort,
information, social pressure), explaining the persistence of political
polarization.

\textbf{Hamiltonian framework (2025).} A novel model provides: (1) complete probabilistic
formulation with memory effects, (2) parameterized interaction potentials for
social hierarchy and network topology, and (3) explicit consensus emergence
mechanisms. Unlike classical opinion dynamics models (DeGroot, bounded confidence),
the quantum model naturally produces oscillatory behavior, tunneling between
opinion clusters, and interference between information pathways.

\textbf{Quantum simulations on quantum hardware (2025).} Models of opinion dynamics have
been implemented on IBM Quantum devices, using superposition for opinion
uncertainty, measurement-induced collapse for opinion commitment, and entanglement
for social correlations.
\subsection{7.2 Information Diffusion via Quantum Walks}
\label{sec:org83ad26e}

Quantum walks --- the quantum analog of classical random walks --- model
information propagation through social networks:

\textbf{Continuous-Time Quantum Walk Information Propagation Model (CTQW-IPM).} Ranks
crucial individuals in social networks for information spread. Overcomes
limitations of classical simulation models (extensive iteration, multiple control
parameters) by using the quantum walk's natural spreading behavior.

\textbf{Asymmetry-resolving quantum walk.} A key innovation addresses the problem that
quantum walk Hamiltonians are Hermitian (symmetric transitions), but social
influence is inherently asymmetric. The model lifts the social network state
into a higher-dimensional space and applies controlled operators to recover
asymmetric influence.

\textbf{Quantum Walk Neural Networks.} Learn a graph-dependent diffusion operation that
can \emph{direct} information flow (unlike classical random walks in GCN/DCNN that
diffuse uniformly). The "coin operators" controlling quantum walk behavior are
learned from data.

\textbf{Graph Quantum Walk Transformer (GQWformer, 2024).} Developed by Zhejiang Lab,
this architecture embeds quantum walks into graph transformers, capturing both
global and local structure. Validated on 5 benchmark datasets across biological,
chemical, and social domains, outperforming 11 baseline models.
\subsection{7.3 Voting Behavior}
\label{sec:org8c0d38a}

Quantum-like models have been applied extensively to voting:

\textbf{Nonseparability of voter preferences.} "Ticket splitting" --- voting for
different parties in executive vs. legislative races --- demonstrates non-separable
preferences that violate classical Kolmogorovian probability through inconsistent
Bayesian updating. The quantum formalism naturally accommodates this via entangled
preference states in the tensor product of the presidential and senatorial
Hilbert spaces.

\textbf{Quantum interference in political preferences.} Election contexts modeled with
quantum probability produce an interference term for contextuality. Some
reconstructed decision eigenvectors produce "hyperbolic interference" (magnitude
exceeding 1), indicating strong non-classicality.

\textbf{Coalition formation (Bagarello, 2015).} Quantum field theory formalism derives
dynamical equations for the evolution of parties' coalition preferences,
incorporating voter behavior impact.

\textbf{Quantum Condorcet Voting and Arrow's Theorem.} Quantum Condorcet Voting has
been used to disprove Arrow's Impossibility Theorem in the quantum setting. For
small readout amplitudes, quantum outcomes track the classical scheme; beyond a
modest noise level, they converge to non-classical distributions that satisfy
conditions Arrow showed to be classically impossible.

\textbf{Density operator models (Dubois).} Density operators relative to candidate
families model voting intention, with proportionality hypotheses between density
matrix coefficients and vote probabilities. Applied to the French 2012
presidential election with promising results.
\subsection{7.4 Social Choice and Contextuality}
\label{sec:org54006d5}

The deepest connection between quantum theory and social science may be through
contextuality:

\textbf{Topological connection (Baryshnikov).} "The Topology of Quantum Theory and Social
Choice" identifies topological singularities encoding the difference between
classical and quantum probability, establishing that a resolution to the social
choice problem is equivalent to a resolution to the violation of unicity in
quantum measurement frameworks.

\textbf{Contextuality-by-Default in economics (2025).} True contextuality has been found
in global investment decisions, demonstrating that quantum-like contextuality
extends beyond cognitive experiments to real-world economic behavior.

\textbf{Intransitivity and contextuality.} The relationship between intransitive
preferences (A ≻ B ≻ C ≻ A) and Type II contextuality (considered sine qua non
of quantum mechanics) has been observed in human decision-making experiments.
\section{8. Categorical and Compositional Approaches}
\label{sec:orgb585e33}

\subsection{8.1 Categorical Quantum Mechanics (CQM)}
\label{sec:org1ac5c23}

Categorical Quantum Mechanics, developed primarily by Bob Coecke, Samson Abramsky,
Chris Heunen, and Jamie Vicary at Oxford, provides the mathematical foundation
for the tensor-aware propagator network design described in the companion document.

The framework axiomatizes quantum processes in terms of:
\begin{itemize}
\item \textbf{Monoidal categories:} Sequential (∘) and parallel (⊗) composition of processes.
\item \textbf{Dagger categories:} Every process f has an adjoint f†, giving time-reversal.
\item \textbf{Compact closed categories:} Duality between objects enables wire-bending in
string diagrams, corresponding to bidirectional constraint propagation.
\item \textbf{Frobenius algebras:} Merge and copy operations at nodes, with the spider theorem
ensuring topology-independence.
\end{itemize}

The graphical calculus (string diagrams) provides a complete reasoning language:
any equation between quantum processes that holds in finite-dimensional Hilbert
spaces can be proved by diagram manipulation alone.
\subsection{8.2 DisCoCat and Quantum NLP}
\label{sec:org5ef3d65}

The most developed social-science-adjacent application of CQM is \textbf{DisCoCat}
(Categorical Compositional Distributional Semantics), introduced by Coecke,
Sadrzadeh, and Clark:

\begin{itemize}
\item \textbf{Key insight:} Pregroup grammars and quantum processes both form rigid monoidal
categories. A monoidal functor maps grammatical structure to vector space
semantics, composing word meanings into sentence meanings via the same algebra
that composes quantum states.
\item \textbf{Verbs as unitary operations:} A transitive verb is a linear map
S ⊗ O → S (where S = subject space, O = object space, S = sentence space).
This is structurally identical to a unitary operator on a tensor product.
\item \textbf{Extensions:} Density matrices for lexical ambiguity (Piedeleu et al.),
convex relations for cognitive modeling (Bolt et al., 2016), and games for
dialogue and Wittgenstein's language games (Hedges \& Lewis, 2018).
\end{itemize}

This connects to Prologos through the quantale-enriched categorical framework:
the same string diagram calculus that describes quantum NLP describes
propagation through TensorCells.
\subsection{8.3 Sheaf-Theoretic Contextuality}
\label{sec:org7f3c0ae}

Abramsky and Brandenburger developed a sheaf-theoretic framework for
contextuality that is \emph{theory-independent} --- it applies to any situation
with locally consistent but globally inconsistent empirical data:

\begin{itemize}
\item Local sections over a measurement context describe the outcomes observable
in that context.
\item The sheaf condition checks whether local sections glue to a global section.
\item Failure of the sheaf condition = contextuality = no single classical
probability model accounts for all the data.
\end{itemize}

This connects directly to Prologos's sheaf-theoretic distributed reasoning
(Beyond Prolog §8.3): the same mathematical structure that checks for
quantum contextuality checks for consistency of distributed knowledge bases.
\subsection{8.4 Connection to Prologos's Categorical Architecture}
\label{sec:orgfec449b}

The categorical structures described in this section align precisely with
Prologos's existing and planned architecture:

\begin{center}
\begin{tabular}{ll}
CQM Structure & Prologos Realization\\
\hline
Monoidal category & Propagator network (sequential + parallel)\\
Dagger & Adjoint/reverse propagation\\
Compact closure & Bidirectional constraints\\
Frobenius algebra & Cell merge/broadcast operations\\
Quantale enrichment & Information-graded morphisms\\
Sheaf condition & Distributed consistency checking\\
Tensor product & TensorCell (joint state spaces)\\
Completely positive maps & DensityCell propagators (measurement)\\
\end{tabular}
\end{center}

The vision: a single computational framework --- the quantale-enriched
propagator network --- that simultaneously supports type inference, session
type checking, relational logic, probabilistic reasoning, and quantum-like
modeling. Each domain is an instantiation of the same enriched-categorical
machinery.
\section{9. Modeling with Tensor-Aware Propagators: Opportunities}
\label{sec:org40f908e}

\subsection{9.1 Decision Modeling}
\label{sec:org2b484dc}

The Yukalov-Sornette QDT framework maps to propagator constraints:

\begin{itemize}
\item Each prospect is a basis element in an \texttt{AmplitudeCell}.
\item The decision-maker's state |ψ⟩ is a superposition of prospects.
\item The attraction factor q(π) emerges from the interference term when
amplitudes are squared: p(π) = |⟨π|ψ⟩|² = f(π) + q(π).
\item The Quarter Law (⟨|q|⟩ = 1/4) becomes a testable propagator constraint.
\item Deliberation (time evolution of preference) is a unitary rotation of |ψ⟩.
\end{itemize}
\subsection{9.2 Game-Theoretic Modeling}
\label{sec:org97f05cd}

The EWL protocol encodes directly:

\begin{itemize}
\item Two-player strategy space = \texttt{TensorCell [Strategy Strategy]}.
\item Entangling gate J(γ) = unitary propagator on the joint TensorCell.
\item Individual strategies = unitary propagators on subsystems.
\item Payoff computation = measurement + classical propagation.
\item Nash equilibrium = propagation to quiescence under payoff constraints.
\end{itemize}

The quantum Prisoner's Dilemma, where entangled strategies produce Pareto-optimal
Nash equilibria, is expressible as a tensor-aware propagator network.
\subsection{9.3 Financial Modeling}
\label{sec:orge089767}

Haven's Schrödinger-Black-Scholes connection:

\begin{itemize}
\item Market states as \texttt{AmplitudeCell} values in superposition.
\item The arbitrage parameter ℏ controls the degree of quantum-like behavior.
\item Unitary evolution encodes fundamental price dynamics.
\item Measurement = trade execution (collapse to definite price).
\item The GKSL equation via \texttt{DensityCell} models market decoherence.
\end{itemize}
\subsection{9.4 Social Network Modeling}
\label{sec:org18816e4}

Quantum walk propagators for opinion dynamics:

\begin{itemize}
\item Each agent's belief is an \texttt{AmplitudeCell} over opinion options.
\item Social influence = unitary rotation toward neighbor's state.
\item Network topology determines propagator wiring.
\item Entanglement between agents = correlated belief updates via \texttt{TensorCell}.
\item Consensus = decoherence to classical mixture via Lindblad operators.
\end{itemize}
\subsection{9.5 Survey and Polling Models}
\label{sec:orgae29d8b}

The QQ equality as a propagator constraint:

\begin{itemize}
\item The Wang-Busemeyer QQ equality is automatically satisfied by any quantum
probability model implemented on AmplitudeCells.
\item Violation detection becomes a constraint-propagation problem.
\item Order effects in survey questions are modeled by non-commutative projections.
\item The propagator network can optimize question ordering to minimize bias.
\end{itemize}
\section{10. Open Problems and Future Directions}
\label{sec:org76bf475}

\subsection{10.1 Empirical Validation Challenges}
\label{sec:orga074828}

While the QQ equality has strong empirical support (70+ national surveys), many
quantum-like models have limited validation beyond explaining known anomalies post
hoc. Prologos's modeling capabilities could help by making it easier to generate
quantitative predictions from quantum-like models and test them against data.
\subsection{10.2 The Parameter Explosion Problem}
\label{sec:org86cc8fc}

Quantum Bayesian networks suffer from exponential growth of parameters with the
number of variables. The similarity heuristic (Moreira \& Wichert) helps, but
principled dimensionality reduction for quantum-like social models remains open.
Tensor network contraction methods may provide a path.
\subsection{10.3 Non-Markovian Extensions}
\label{sec:orgde5e3aa}

The GKSL master equation assumes Markovian (memoryless) dynamics, but human
cognition has strong memory effects. Non-Markovian quantum dynamics
(Nakajima-Zwanzig equations, time-convolutionless approaches) are an active
frontier. Prologos's persistent (immutable) propagator network naturally
preserves historical states, which could be exploited for memory kernels.
\subsection{10.4 Connecting Quantum-Like to Actual Quantum Computing}
\label{sec:org8fbfdfd}

As quantum hardware matures (IBM Heron, IonQ), practical applications in finance
are demonstrating near-parity with classical methods. The question arises: can
quantum-like cognitive/social models be run on actual quantum hardware for
speedup? The answer is nuanced --- the models use quantum math but typically
involve classical simulation of small quantum systems. Quantum advantage would
emerge for large-scale social simulations (many agents, entangled beliefs).
\subsection{10.5 Contextuality as the Key Resource}
\label{sec:org03f4c3f}

The CbD framework (Dzhafarov \& Kujala) provides a theory-independent way to
detect and measure contextuality in behavioral data. Contextuality may be the
key "resource" that quantum-like models provide --- analogous to how entanglement
is a resource in quantum computing. Understanding when and why social/economic
systems are contextual would guide where quantum-like models are most valuable.
\subsection{10.6 Cross-Domain Compositional Models}
\label{sec:org2f4401a}

The categorical framework (CQM, DisCoCat) suggests the possibility of
\emph{compositional} quantum-like models that span domains: a model of voter cognition
(quantum decision theory) composed with a model of social influence (quantum
walks) composed with a model of media dynamics (open quantum systems). Prologos's
propagator network, with its natural compositional structure, is well-suited to
this vision.
\section{11. References}
\label{sec:org07ec632}

\subsection{Quantum Decision Theory}
\label{sec:org2be310c}
\begin{itemize}
\item Yukalov, V.I. \& Sornette, D. (2010). Mathematical structure of quantum decision theory. \emph{Advances in Complex Systems} 13, 659-698.
\item Yukalov, V.I. \& Sornette, D. (2011). Decision theory with prospect interference and entanglement. \emph{Theory and Decision} 70, 283-328.
\item Yukalov, V.I. \& Sornette, D. (2014). Conditions for quantum interference in cognitive sciences. \emph{Topics in Cognitive Science} 6, 79-90.
\item Yukalov, V.I. \& Sornette, D. (2016). Quantum probabilities as behavioral probabilities. \emph{Entropy} 19, 112.
\end{itemize}
\subsection{Quantum Cognition}
\label{sec:org36473e1}
\begin{itemize}
\item Busemeyer, J.R. \& Bruza, P.D. (2012). \emph{Quantum Models of Cognition and Decision}. Cambridge UP.
\item Pothos, E.M. \& Busemeyer, J.R. (2013). Can quantum probability provide a new direction for cognitive modeling? \emph{Behavioral and Brain Sciences} 36, 255-274.
\item Wang, Z. \& Busemeyer, J.R. (2013). A quantum question order model supported by empirical tests of an a priori and target prediction. \emph{Topics in Cognitive Science} 5, 689-710.
\item Wang, Z. et al. (2014). Context effects produced by question orders reveal quantum nature of human judgments. \emph{PNAS} 111, 9431-9436.
\end{itemize}
\subsection{Quantum Game Theory}
\label{sec:org849fd9c}
\begin{itemize}
\item Eisert, J., Wilkens, M. \& Lewenstein, M. (1999). Quantum games and quantum strategies. \emph{Physical Review Letters} 83, 3077.
\item Benjamin, S.C. \& Hayden, P.M. (2001). Multiplayer quantum games. \emph{Physical Review A} 64, 030301.
\item Bagarello, F. (2015). Quantum dynamics for classical systems. \emph{Springer}.
\end{itemize}
\subsection{Quantum Finance}
\label{sec:orgddb1e43}
\begin{itemize}
\item Haven, E. \& Khrennikov, A. (2013). \emph{Quantum Social Science}. Cambridge UP.
\item Haven, E. (2002). A discussion on embedding the Black-Scholes option pricing model in a quantum physics setting. \emph{Physica A} 304, 507-524.
\item Rao, A., Yang, S. \& Zakerinia, S. (2025). Using quantum game theory to model competition. \emph{SSRN}.
\end{itemize}
\subsection{Open Quantum Systems in Cognition}
\label{sec:org0e8611a}
\begin{itemize}
\item Asano, M. et al. (2011). Quantum-like model of brain's functioning. \emph{Journal of Theoretical Biology} 281, 56-64.
\item Asano, M. et al. (2011). Quantum-like dynamics of decision-making. \emph{Physica A} 391, 2083-2099.
\item Khrennikov, A. (2023). Open systems, quantum probability, and logic for quantum-like modeling. \emph{Entropy} 25, 886.
\end{itemize}
\subsection{Contextuality}
\label{sec:orge5a72cb}
\begin{itemize}
\item Dzhafarov, E.N. \& Kujala, J.V. (2016). Context-content systems of random variables. \emph{Philosophical Transactions A} 374, 20150235.
\item Cervantes, V.H. \& Dzhafarov, E.N. (2018). Snow queen is evil and beautiful. \emph{Decision} 5, 193-214.
\item Baryshnikov, Y. (2022). The topology of quantum theory and social choice. \emph{Quantum Reports} 4, 14.
\end{itemize}
\subsection{Quantum Walks and Social Networks}
\label{sec:org14df466}
\begin{itemize}
\item Li, X. et al. (2022). CTQW information propagation model. \emph{Neural Computing and Applications}.
\item Dernbach, S. et al. (2019). Quantum walk neural networks. \emph{Applied Network Science}.
\end{itemize}
\subsection{Categorical and Compositional}
\label{sec:org14d4051}
\begin{itemize}
\item Coecke, B., Sadrzadeh, M. \& Clark, S. (2010). Mathematical foundations for a compositional distributional model of meaning. \emph{Linguistic Analysis} 36, 345-384.
\item Abramsky, S. \& Brandenburger, A. (2011). The sheaf-theoretic structure of non-locality and contextuality. \emph{New Journal of Physics} 13, 113036.
\end{itemize}
\subsection{Prologos Internal Documents}
\label{sec:org453a76e}
\begin{itemize}
\item Beyond Prolog: Frontiers for the Prologos Relational Language (2026-03-16).
\item Tensor-Aware Propagator Networks with Unitary Evolution (companion, 2026-03-16).
\item Logic Engine Design: Comprehensive Implementation Guide (2026-02-24).
\item Towards a General Logic Engine on Propagators (2026-02-24).
\end{itemize}
\end{document}
